\subsection{IC in DIP Package (AT28C64U / General-Purpose)}
\markright{IC in DIP Package}

\subsubsection*{Definition}
An integrated circuit (IC) in a dual in-line package (DIP) is a semiconductor device
in which multiple transistors, resistors, and other elements are fabricated on a single
silicon die and enclosed in a rectangular plastic or ceramic body with two parallel rows
of pins on 2.54\,mm pitch \cite{sedra2020}. The DIP format is compatible with both
breadboard and veroboard prototyping. The 28-pin DIP device in this laboratory
(tentatively identified as an AT28C64U type EEPROM) is a representative example of how
digital memory and logic ICs are packaged for laboratory use.

\labimage{images/dip28_ic.jpg}%
  {A 28-pin DIP integrated circuit as used in the EEE 2188 laboratory.
   The notch at one end identifies pin 1; pins are numbered
   counter-clockwise from pin 1.}{dip28_lab}

\subsubsection*{Specifications}
DIP packages are available from 6 to 64 pins. The 28-pin DIP has an overall body width
of 0.6\,inch (15.24\,mm) and an overall pin-to-pin span of 0.1\,inch pitch. IC
specifications vary entirely by device type---digital, analog, memory, or mixed-signal.

\subsubsection*{Purpose of Use}
DIP ICs are inserted into breadboard or veroboard for prototyping digital logic, memory,
microcontroller, and analog IC circuits. The DIP format makes pin identification and
probing straightforward.

\subsubsection*{Applications}
DIP-packaged ICs include logic gates (74xx TTL/CMOS series), op-amps (LM741, LM324),
comparators (LM339), timer ICs (NE555), SRAM and EEPROM memory chips, and
microcontrollers (PIC, ATmega) \cite{sedra2020}.
